% sections/summary.tex -- shared abstract/summary.
% \input by proposal-rebecca.tex (Summary) and formal-proposal-draft.tex
% (non-technical Summary required by the Dissertation Regulations).
% Edit here once; both documents change. No sectioning commands in this file.
I am an artistic practitioner and researcher working through
research-creation. Trained in architecture and landscape architecture, I build
interactive 3D environments and multiplayer games and study them with
participants. This project asks how generative artificial intelligence
(GenAI), meaning the models, interfaces and platforms that operate them,
reorganizes interpretive authority and collective authorship in
collaborative artistic practice. Three subquestions follow: how selected
tools, defaults and facilitation arrangements amplify, standardize, misread
or pass over participants' culturally situated contributions; how
collaborators negotiate uptake, refusal, credit and responsibility among
themselves, the group, the artist-facilitator, models, platforms and
institutions; and how artistic methods and a public artwork can make these
usually hidden distributions perceptible, debatable and revisable. As
platforms mediate more cultural production, they help decide whose meanings
become legible and who answers for shared work. The design begins with my own studio inquiry, followed by
consultation with four to six practitioners, workshops with 18--24 adult
interdisciplinary artists in three cohorts with participant observation,
semi-structured focus groups, and the co-creation of a room-scale
collaborative installation proposed for exhibition at York University. It
would contribute a communication-and-culture account of how GenAI
redistributes interpretive authority, a research-creation method for
studying it, and an artwork that keeps those negotiations visible.
